%% AIQxQIA 2026 full paper (>4 pages + references), CEUR-WS ceurart class,
%% single-blind (authors named). Submission via EasyChair.
\documentclass{ceurart}

\usepackage{amsmath,amssymb}
\usepackage{booktabs}

\newcommand{\GF}{\mathbb{F}_2}
\newcommand{\HX}{H_X}
\newcommand{\HZ}{H_Z}
\newcommand{\ranktwo}{\operatorname{rank}_{\GF}}
\newcommand{\wt}{\operatorname{wt}}
\newcommand{\eff}{kd^2/n}

\begin{document}

\copyrightyear{2026}
\copyrightclause{Copyright for this paper by its authors. Use permitted under
    Creative Commons License Attribution 4.0 International (CC BY 4.0).}
\conference{AIQxQIA 2026: International Workshop on AI for Quantum and Quantum
    for AI, 2026}

\title{The QEC Challenge: an adversarially verified leaderboard for AI-driven
    discovery of quantum LDPC codes}

%% Authors alphabetical by surname. TODO before submission: add ORCIDs and
%% the co-authors' emails.
\author[1]{Bradley A. Chase}
\author[1]{Sebastian Hassinger}
\author[1]{Farrokh Labib}
\author[1]{Seun Omonije}
\author[1]{Mathys Rennela}
\author[1]{Vincent Russo}[email=vincentrusso1@gmail.com]
\author[1,2]{William J. Zeng}
\address[1]{Unitary Foundation}
\address[2]{Quantonation}

\begin{abstract}
    AI systems now produce candidate quantum low-density parity-check (qLDPC)
    codes faster than human referees can check them, and the characteristic
    failure mode of machine-generated codes is an overstated distance. The QEC
    Challenge is a public leaderboard built around that fact: a submission is a
    single machine-readable description of a CSS code, and an automated
    pipeline recomputes every cheap property from scratch, verifies a
    self-certifying distance witness, and then actively attacks the distance
    claim with independent randomized searches before the code is accepted.
    Three design principles follow. Everything inexpensive is trustless and
    mandatory; distance confidence is layered, from a witnessed upper bound
    through corroboration to exact certification by integer programming; and
    the ranking is a Pareto frontier inside computed constraint cells rather
    than a single scalar. Because nothing a submission contains is trusted,
    the leaderboard doubles as the verification harness of LLM- and
    agent-driven code-discovery loops, and several of its current records were
    found by such loops. We describe the design, report on the first months of
    operation (124 verified codes, 29 with exactly certified distance), and
    present the empirical lessons the board has produced about how deeply a
    heuristic distance claim must be attacked before it deserves confidence.
\end{abstract}

\begin{keywords}
    quantum error correction \sep qLDPC codes \sep benchmark \sep
    leaderboard \sep verification \sep large language models \sep
    automated discovery
\end{keywords}

\maketitle

\section{Introduction}

Quantum low-density parity-check (qLDPC) codes are the leading route to
fault-tolerant quantum memories with constant or near-constant overhead:
asymptotically good families exist~\cite{panteleev2022asymptotically}, and
concrete constructions such as the bivariate-bicycle codes of Bravyi et
al.~\cite{bravyi2024high} encode an order of magnitude more logical qubits
than surface-code patches of comparable distance. Finding good codes at
practical block lengths, however, remains a search problem over a large
discrete space, and it is exactly the kind of search that machine methods have
begun to win. LLM-guided evolutionary search~\cite{cruzbenito2026evolutionary},
reinforcement learning~\cite{he2025reinforcement}, and LLM-driven concept
evolution over non-abelian code families~\cite{liu2026concept} all generate
candidate codes at a rate no human refereeing process can keep up with.

The bottleneck this creates is verification, and one property in particular.
Computing the distance of a stabilizer code is
NP-hard~\cite{kapshikar2023hardness}, and the hardness is asymmetric:
exhibiting a single low-weight logical operator proves an upper bound, while
certifying that no lighter one exists is the intractable direction. Heuristic
distance estimates, the kind a search loop produces internally, deflate
routinely under deeper search. A leaderboard that accepted claimed distances
at face value would fill with codes whose headline numbers are artifacts of
under-powered estimation, and a machine-driven search has every incentive,
even an unintentional one, to overclaim.

The QEC Challenge\footnote{\url{https://github.com/unitaryfoundation/qldpc-challenge}}
is a public, continuously verified leaderboard designed around that asymmetry.
This paper describes its design and reports on its first months of operation.
The system faces both directions of the AI and quantum intersection: it is
infrastructure that makes AI-driven quantum code discovery auditable, and its
accumulated history of refuted and confirmed claims is a dataset about the
reliability of heuristic verification, a question AI methods for quantum error
correction cannot avoid.

\section{Design principles}

Three principles shape the system.

\paragraph{Everything cheap is trustless and mandatory.} A submission
describes one CSS code by its two parity-check matrices $\HX$ and $\HZ$,
given as sparse supports, together with claimed parameters $[[n,k,d]]$,
optional geometric layout data, and provenance. The verifier recomputes every
polynomial-time property from scratch over $\GF$: the qubit count implied by
the checks, the logical dimension
\begin{equation}
    k \;=\; n-\ranktwo(\HX)-\ranktwo(\HZ),
\end{equation}
CSS commutation $\HX\HZ^{\top}=0$, the maximum check weight, and the locality
class implied by the layout. A wrong claimed $k$ or a non-commuting
``stabilizer'' is rejected in milliseconds. None of these facts are taken from
the submission; the submission merely has to agree with them.

\paragraph{Distance confidence is layered by how much is proven.} The
distance claim is split into an upper bound the submitter proves and a lower
bound the submitter is not asked for. Every submission must include, for each
Pauli side, a \emph{witness}: the support of an explicit logical operator of
the claimed weight. The verifier checks that the witness lies in the kernel
of the opposite-side checks and outside the row space of its own side, which
certifies $d \le \text{value}$ with no trust required. Everything beyond the
upper bound is earned in tiers described in
Section~\ref{sec:distance}: an adversarial refutation gate that must fail to
beat the claim, an optional corroboration tier recording how hard the claim
has been attacked, and an exact tier in which an integer program proves no
lighter logical exists.

\paragraph{Ranking by Pareto dominance.} A quantum code trades physical
qubits $n$, logical qubits $k$, distance $d$, check weight $w$, and geometric
locality against one another. Collapsing these into one scalar invites gaming
that scalar. The board instead stratifies codes into cells by two computed
constraint axes and ranks each cell by Pareto dominance, as described next.

\section{Cells and Pareto frontiers}

Track membership is computed by the verifier from the parity checks and the
layout; it cannot be claimed. The two axes are a locality class (single-layer
2D-local, bilayer 2D-local in the flip-chip regime, or unrestricted; the
class is earned from submitted coordinates whose interaction radius the
verifier measures, and a missing or dishonest layout simply lands the code in
unrestricted) and a check-weight class ($w\le 4$, $w\le 6$, $w\le 8$, or any
weight, from the maximum row weight of $\HX$ and $\HZ$). The classes nest: a
single-layer weight-4 code competes in every looser cell as well.

Within a cell, a code is a record if no other code in the cell dominates it,
meaning no other code has $n'\le n$, $k'\ge k$, $d'\ge d$, and $w'\le w$ with
at least one strict inequality. There can be many co-leaders, and that is
deliberate: a cell is a frontier, not a podium. Construction-family labels
(bivariate bicycle, generalized bicycle and its non-abelian two-block
group-algebra generalization~\cite{lin2024quantum}, coset, tile, topological)
cannot be recovered from the matrices, so they are filterable tags with no
ranking effect.

The board is seeded with literature baselines, including the planar
open-boundary codes of Liang, Eberhardt, and Chen~\cite{liang2025planar}, the
gross code~\cite{bravyi2024high}, generalized-bicycle
codes~\cite{panteleev2021degenerate,wang2022distance}, tile
codes~\cite{steffan2025tile}, two-block coset codes~\cite{aydin2026breaking},
and the surface, toric, and Steane
codes~\cite{kitaev2003fault,steane1996error}, so a new entry is measured
against published codes rather than an empty table. The board complements the
Error Correction Zoo~\cite{eczoo}: the Zoo catalogs code families and their
known parameters, while the board stores concrete parity-check instances whose
claims are machine-checked and then kept under attack.

\section{Figures of merit}
\label{sec:metrics}

The efficiency figure $\eff$ is the standard headline number in the qLDPC
literature, usually motivated by the Bravyi--Poulin--Terhal (BPT)
bound~\cite{bravyi2010tradeoffs}: a geometrically 2D-local code obeys
$kd^2 \le c\,n$ with $c$ independent of $n$. Operating the board forced us to
be precise about what this number means, because most codes on it are
expander-like and not 2D-local, so the BPT bound does not bind them at all,
and for asymptotically good families $\eff$ grows like $n^2$ without bound.

The reading that survives is baseline-relative: $k$ rotated surface-code
patches of distance $d$ cost $kd^2$ data qubits, so $\eff$ is the
qubit-savings factor against building the same logical count and distance out
of the incumbent architecture. That reading is legitimate on any hardware,
needs no geometry theorem, and explains the metric's behavior; it also makes
clear that the comparison concedes long-range connectivity for free. The
board therefore reports $\eff$ as a sortable, incumbent-relative column and
never lets it collapse a Pareto frontier.

For codes that do certify a geometric layout, the board reports a second,
ceiling-relative score obtained from the BPT tradeoff with its range and
dimension dependence made explicit,
\begin{equation}
    f \;=\; \frac{16\, k\, d^{2/(D-1)}}{w_r^{\,2+2/(D-1)}\, n},
\end{equation}
where $D$ is the grid dimension and $w_r$ the interaction range measured from
the submitted layout. The exponents are provably optimal (tiling a good qLDPC
code~\cite{panteleev2022asymptotically} into range-$w_r$ boxes saturates
them), and the constant is fixed so that the rotated surface code scores
exactly $f=1$ at every distance; any 2D-local code with $f>1$ would beat the
surface code at its own geometric game. The best 2D-local code currently on
the board reaches $f=0.907$, and whether $f>1$ is achievable by a 2D-local
family is an open target the board makes concrete. The two scores answer
different questions, savings against the incumbent versus saturation of a
proven ceiling, and the board keeps them separate by construction.

\section{Distance verification: witness, attack, certificate}
\label{sec:distance}

\paragraph{Trustless upper bound.} The submitted witness is validated
arithmetically, as described above. This tier requires no trust in the
submitter, which is what lets anonymous or automated contributors climb the
board instantly.

\paragraph{The adversarial gate.} Every new or changed code then faces
independent bounded searches for a logical operator \emph{lighter} than the
claim: a random-information-set search in the spirit of
QDistRnd~\cite{pryadko2023qdistrnd} and a syndrome-decoder search built on
belief propagation with ordered-statistics
post-processing~\cite{roffe2020decoding}, two different search mechanisms.
Budgets are adaptive. Trial counts scale with block length, and a code that
would enter a cell's Pareto frontier faces a much deeper multi-seed battery
than a dominated one, so scrutiny concentrates exactly where gaming would
pay. Any hit is a sound refutation, since it exhibits a checkable
counterexample, and rejects the submission. Seeds are fresh on every run, so
an overclaim cannot be tuned to survive one fixed search, and a weekly sweep
re-attacks the entire board to accumulate coverage over time.

\paragraph{Exact certification.} A separate, maintainer-run tier upgrades an
upper bound to an exact distance by integer programming: for each logical
generator $t$ of the relevant type, a bounded solver decides feasibility of
\begin{equation}
    H v \equiv 0 \pmod 2, \qquad
    \langle v, t\rangle \equiv 1 \pmod 2, \qquad
    \wt(v) \le d-1,
\end{equation}
with the parity constraints linearized by integer slacks. Infeasibility for
every generator on both sides proves $d$ exact; a feasible point is a
refutation; a solver timeout leaves the claim an upper bound.
Nothing is ever silently upgraded, and every tier is backed by an artifact in
the repository: the witness inside the submission, a heuristic certificate
recording the failed attack, or an exact certificate. Of the 124 codes
currently on the board, 29 carry exactly certified distances; the integer
programs currently stop scaling near $n\approx 300$.

\section{Trust architecture for machine submissions}

The pipeline assumes every submission may come from an automated system and
that nothing it contains can be trusted. Four mechanisms carry that
assumption.

First, scope separation: a pull request that adds code data may not touch the
verifier, the schema, the workflows, or the site generator, and for code
submissions the entire judgment runs from the base branch, so a submission
cannot alter the rules it is judged by. Second, a pinned validation stack:
the trusted closure of the verifier, including the refutation engine, the
$\GF$ core, and the integrity checker itself, is pinned by a cryptographic
manifest, and any drift fails the build, making changes to the rules a
deliberate, reviewable act. Third, authorship binding: the submitting account
must be listed as an author of the code it adds. Fourth, duplicate detection:
row-echelon forms of $(\HX,\HZ)$ pin the exact stabilizer group, and a
permutation-invariant Weisfeiler--Leman signature of the Tanner graph flags
likely equivalents for human review.

Two further conventions target AI-driven contributors specifically.
Submissions carry an optional, self-reported model attribution naming the
system that produced the code to a specific version, so machine-discovered
entries are labeled as such on the board. And the same validator that CI runs
is available locally, so an autonomous search agent can adopt the discipline
that no candidate counts as a find until the pinned validator passes it,
reaching the same verdict CI will reach before a pull request is ever opened.

Finally, the board enforces a block-length cap as a verification-budget rule.
The adaptive gate runs under a fixed wall-clock budget, so its trials per
qubit thin out as $n$ grows, and exact certification scales worse; above the
cap the pipeline can no longer stand behind a distance claim at all. The cap
states what the machinery can vouch for today and is raised as the tooling
improves, keeping the board a finite-length benchmark rather than
letting scale substitute for evidence.

\section{Operational lessons}
\label{sec:lessons}

At the time of writing the board holds 124 verified codes across the
bivariate-bicycle, generalized-bicycle, two-block group-algebra, coset,
lifted-product, tile, and topological families, a majority contributed
through the challenge rather than seeded, with the best incumbent-relative
efficiency at $\eff = 189.2$. Several records are machine-discovered and
attributed to the producing model, including the current efficiency leader, a
$[[390,82,30]]$ generalized-bicycle code found by an LLM-driven search over
divisor structures of $x^{195}-1$, a $[[336,20,20]]$ two-block code on the
Klein-quartic group $\mathrm{PSL}(2,7)$, and a $[[240,13,15]]$ code on the
symmetric group $S_5$, the board's first odd-$k$ entry (odd $k$ is impossible
for abelian two-block constructions).

The board's history has begun to answer a question that any AI-for-QEC
pipeline must face: how deep must a randomized refutation search run before
an upper-bound claim deserves confidence? The record so far:

\begin{itemize}
    \item In one agent-driven sweep, five of seven staged frontier candidates
    were refuted by the deep gate before a pull request was opened; the
    shallow estimates that had promoted them were optimistic in every case.
    \item Around the current efficiency record, eleven independent candidates
    from the same construction family read distances of $31$ to $34$ under
    $8\times 10^6$ randomized trials, an apparent improvement on the record.
    Under $8\times 10^7$ trials per seed across three seeds, all eleven
    collapsed to $d\le 30$, exactly the record they appeared to beat. The
    factor-of-ten deeper search changed the verdict in every single case.
    \item A submitted $[[684,8]]$ code claimed $d=100$; the gate's two
    independent mechanisms found logicals of weight $97$ and then $85$,
    and the entry became $d\le 85$. The refutation arrived in
    continuous integration, minutes after submission.
\end{itemize}

The general pattern is that heuristic distance estimates at $10^6$ to $10^7$
trials are systematically inflated at block lengths in the low hundreds, and
that independent search mechanisms (information-set sampling versus
syndrome decoding) catch overclaims that either alone would miss. For
machine-discovery loops, whose internal fitness signal is exactly such a
heuristic estimate, this argues for a verification budget comparable to the
search budget, and for treating any frontier-adjacent estimate as provisional
until it has survived a multi-seed attack an order of magnitude deeper than
the search that produced it.

\section{Open questions}

The board keeps several questions concrete and live. Can planar
open-boundary codes close the efficiency gap to periodic and unrestricted
records, or is part of that gap intrinsic to boundaries? Can any 2D-local
family exceed the surface code's ceiling-relative score, $f>1$? How far does
the non-abelian two-block design space reach, given that groups outside the
systematically enumerated ranges keep yielding records? Can exact
certification scale past $n\approx 300$, and can refutation-depth
requirements be predicted from $(n,k,w)$ rather than learned by collapse? The
challenge is open, the pipeline and its history are public, and both human
and automated contributors are welcome to attack the frontiers, and the
claims, directly. The design is meant to scale with participation:
verification is automated and per-submission, so widening the contributor
pool adds search volume without diluting the board's guarantees, and every
accepted or refuted claim enriches the public record either way. Whether an
open, adversarially verified frontier actually speeds up code discovery is
itself an experiment this challenge runs; the early signs, records set by
independent contributors and search loops within weeks of launch, are
consistent with it, though the sample is small.

%% ceurart's default style (elsarticle-num-names) is vendored alongside.
\bibliography{refs}

\end{document}
